\section*{الفصل \ref{ch:g10:proba} --- الاحتمالات والمعاينة}

\begin{solution}{exo:g10:proba:1}
\omterm{def:g10:proba:distribution}{احتمال} كل إمكانية من الإمكانيات $6$ هو $\frac16$.
``الحصول على $6$'': $\frac16$.
``عدد فردي'' $= \{1, 3, 5\}$: $\frac36 = \frac12$.
``$4$ على الأكثر'' $= \{1,2,3,4\}$: $\frac46 = \frac23$.
``الحصول على $7$'': \omterm{def:g10:proba:events}{حادثة} مستحيلة، واحتمالها $0$.
\end{solution}

\begin{solution}{exo:g10:proba:2}
توجد $10$ كرات متساوية \omterm{def:g10:proba:distribution}{الاحتمال}.
$\P(\text{أخضر}) = \frac{5}{10} = \frac12$.
$\P(\text{ليست سوداء}) = 1 - \frac{2}{10} = \frac{8}{10} = \frac45$.
والحادثتان ``خضراء'' و``صفراء'' متنافيتان:
$\P(\text{أخضر أو أصفر}) = \frac{5}{10} + \frac{3}{10} = \frac{8}{10}
= \frac45$.
\end{solution}

\begin{solution}{exo:g10:proba:3}
لتكن $S$: ``يدرس الإسبانية''، و $G$: ``يدرس الألمانية''. عندئذٍ
$\P(S) = \frac{18}{30}$ و $\P(G) = \frac{10}{30}$ و
$\P(S \cap G) = \frac{4}{30}$. وبقاعدة الجمع:
\[
\P(S \cup G) = \frac{18}{30} + \frac{10}{30} - \frac{4}{30}
= \frac{24}{30} = \frac45 .
\]
وبالمتممة: $\P(\text{لا هذه ولا تلك}) = 1 - \frac45 = \frac15$.
\end{solution}

\begin{solution}{exo:g10:proba:4}
\omterm{met:g10:proba:tree}{للشجرة} أربعة مسارات: صورة-صورة، وصورة-كتابة، وكتابة-صورة، وكتابة-كتابة، \omterm{def:g10:proba:distribution}{واحتمال} كل منها
$\frac12 \times \frac12 = \frac14$.
صورتان: $\frac14$. وصورة واحدة بالضبط (مساران): $\frac24 = \frac12$.
وصورة واحدة على الأقل: $1 - \P(\text{TT}) = 1 - \frac14 = \frac34$.
\end{solution}

\begin{solution}{exo:g10:proba:5}
السحب الأول: حمراء \omterm{def:g10:proba:distribution}{باحتمال} $\frac{4}{10}$، وبيضاء \omterm{def:g10:proba:distribution}{باحتمال} $\frac{6}{10}$.
والسحب الثاني \emph{دون} إرجاع: بعد حمراء، تبقى $3$ حمراء و $6$
بيضاء من بين $9$؛ وبعد بيضاء، تبقى $4$ حمراء و $5$ بيضاء من بين $9$.

حمراوان: $\frac{4}{10} \times \frac39 = \frac{12}{90} = \frac{2}{15}$.

ومختلفتا اللون، أي المساران حمراء-بيضاء وبيضاء-حمراء:
\[
\frac{4}{10} \times \frac69 + \frac{6}{10} \times \frac49
= \frac{24}{90} + \frac{24}{90} = \frac{48}{90} = \frac{8}{15}.
\]
\end{solution}

\begin{solution}{exo:g10:proba:6}
\emph{1.} يُظهر كل حجر نرد، على نحو مستقل، كل وجه \omterm{def:g10:proba:distribution}{باحتمال}
$\frac16$، إذن الأزواج المرتبة $36$ متساوية \omterm{def:g10:proba:distribution}{الاحتمال}؛ ويتحدد المجموع
بالزوج.

\emph{2.} المجموع $7$: الأزواج $(1,6), (2,5), (3,4), (4,3), (5,2), (6,1)$،
إذن $\P = \frac{6}{36} = \frac16$. والمجموع $12$: الزوج $(6,6)$ وحده، إذن
$\P = \frac{1}{36}$. والمجموع $4$ على الأكثر: الأزواج التي مجموعها $2$ أو $3$ أو $4$:
$(1,1)$؛ و $(1,2), (2,1)$؛ و $(1,3), (2,2), (3,1)$ --- أي ستة أزواج، إذن
$\P = \frac{6}{36} = \frac16$.

\emph{3.} ينمو عدد الأزواج المحققة لكل مجموع من $1$ (للمجموع $2$)
إلى $6$ (للمجموع $7$) ثم يتناقص: فالمجموع الأكثر \omterm{def:g10:proba:distribution}{احتمالًا} هو $7$.
\end{solution}

\begin{solution}{exo:g10:proba:7}
يُجاب عن كل سؤال إجابة صحيحة \omterm{def:g10:proba:distribution}{باحتمال} $\frac14$، على نحو
مستقل. \omterm{met:g10:proba:tree}{والشجرة} ذات مستويين:

كلاهما صحيح: $\frac14 \times \frac14 = \frac{1}{16}$.

وواحد بالضبط صحيح، أي المساران صحيح-خطأ وخطأ-صحيح:
$\frac14 \times \frac34 + \frac34 \times \frac14 = \frac{6}{16}
= \frac38$.

ولا واحد صحيح: $\frac34 \times \frac34 = \frac{9}{16}$.

تحقق: $\frac{1}{16} + \frac{6}{16} + \frac{9}{16} = 1$.
\end{solution}

\begin{solution}{exo:g10:proba:8}
\emph{1.} \omterm{def:g10:stats:series}{التواتر} المرصود: $\frac{396}{900} = 0.44$. ومن أجل $n = 900$:
$\frac{1}{\sqrt{900}} = \frac{1}{30} \approx 0.033$، إذن \omterm{def:g10:numbers:interval}{فترة} التذبذب
هي نحو $\intcc{0.467}{0.533}$.

\emph{2.} تقع القيمة المرصودة $0.44$ خارج \omterm{def:g10:numbers:interval}{الفترة}: فعينات من $900$
ناخب مأخوذة من مجتمع نسبة الرضا فيه $50\%$ نادرًا ما تشرد إلى هذا الحد.
إذن يلقي الاستطلاع شكًّا جديًّا على الادعاء $50\%$.
\end{solution}

\begin{solution}{exo:g10:proba:9}
\emph{1.} $\P(\text{ربح } 6) = \frac16$ (الحصول على $6$)، و
$\P(\text{ربح } 3) = \frac16$ (الحصول على $5$)، و
$\P(\text{ربح } 0) = \frac46 = \frac23$.

\emph{2.} على \omterm{def:g10:stats:quartiles}{مدى} $600$ لعبة، نتوقع نحو $100$ ستة و $100$ خمسة و $400$
رمية أخرى. والمكاسب: نحو $100 \times 6 + 100 \times 3 = 900$. والتكلفة:
$600 \times 2 = 1200$. والنتيجة الصافية المنتظرة: $900 - 1200 = -300$، أي
خسارة قدرها نحو $0.5$ لكل لعبة في \omterm{def:g10:stats:mean}{المتوسط} --- فاللعبة لا تستحق
أن تُلعب.
\end{solution}

\begin{solution}{pb:g10:proba:1}
\textbf{1.} $P(\text{ملك أو كوبة}) = \frac{4}{52} +
\frac{13}{52} - \frac{1}{52} = \frac{16}{52} = \frac{4}{13}$:
فملك الكوبة يقع في الحادثتين ويجب أن لا يُحسب
مرتين --- ومن هنا الطرح
(\cref{thm:g10:proba:addition}).

\textbf{2.} $P(\text{لا صورة}) = \left(\frac12\right)^3 =
\frac18$، إذن $P(\text{واحدة على الأقل}) = \frac78$.

\textbf{3.} $P(RR) = \frac35 \times \frac24 = \frac{3}{10}$؛
$P(\text{واحدة من كل}) = \frac35 \times \frac24 + \frac25
\times \frac34 = \frac{6}{20} + \frac{6}{20} = \frac35$.

\textbf{4.} $P(BB) = \frac25 \times \frac14 = \frac1{10}$؛ و
$\frac{3}{10} + \frac{6}{10} + \frac{1}{10} = 1$: فمجموع أوراق
\omterm{met:g10:proba:tree}{الشجرة} يساوي $1$ دائمًا.

\textbf{5.} تخضع \omterm{def:g10:proba:operations}{الحادثتان المتنافيتان} للعلاقة $P(A \cup B) = P(A) +
P(B) = 1.1 > 1$: وهذا مستحيل --- فقد اختُلق بعض \omterm{def:g10:proba:distribution}{الاحتمال}.

\textbf{6.} (يجيب معظم الناس، ومنهم كتّاب الرسائل الألف،
بقولهم ``لا يهم، النصف بالنصف''. احتفظ بجوابك من أجل
السؤال 8.)

\textbf{7.} السيارة خلف الباب 1 (\omterm{def:g10:proba:distribution}{باحتمال} $\frac13$): يفتح المقدّم
أيًّا من بابَي العنزتين؛ والتبديل يترك السيارة: \emph{خسارة}.
السيارة خلف الباب 2 ($\frac13$): \emph{يجب} أن يفتح المقدّم الباب 3؛
والتبديل يحط على الباب 2: \emph{فوز}. السيارة خلف الباب 3:
بالتناظر: \emph{فوز}. والمجموع:
$P(\text{التبديل يفوز}) = \frac13 + \frac13 = \frac23$.

\textbf{8.} يكون اختيارك الأول خاطئًا \omterm{def:g10:proba:distribution}{باحتمال}
$\frac23$ --- وعندئذٍ بالضبط، يترك كشف المقدّم المُلزَم
السيارةَ خلف الباب الباقي: فالتبديل يحوّل
كل خطأ أول إلى فوز. إذن $\frac23$، مهما قال
الحدس.

\textbf{9.} ستة سيناريوهات متساوية \omterm{def:g10:proba:distribution}{الاحتمال}: السيارة خلف $1$ أو $2$
أو $3$، والمقدّم يفتح الباب $2$ أو $3$ (ولا يفتح بابك أبدًا). ويكشف الفتح
السيارة في اثنين منها (سيارة عند 2 وفتح 2، وسيارة عند 3 وفتح 3):
فيُستبعدان. وبين السيناريوهات الأربعة الباقية --- سيارة عند 1 وفتح 2،
وسيارة عند 1 وفتح 3، وسيارة عند 2 وفتح 3، وسيارة عند 3 وفتح 2 --- يفوز التبديل
في الأخيرين: $\frac24 = \frac12$. فمع مقدّم جاهل، يكون
``النصف بالنصف'' \emph{صحيحًا}: إذ تعيش المفارقة كلها في
معرفة المقدّم، التي ترشّح السيناريوهات ترشيحًا غير متكافئ.

\textbf{10.} يحمل اختيارك السيارة \omterm{def:g10:proba:distribution}{باحتمال}
$\frac{1}{100}$؛ وتصبّ كشوف المقدّم العليم $98$
الاحتمالَ الباقي $\frac{99}{100}$ على الباب المغلق الوحيد. فبدّل،
واربح $99$ مرة من كل $100$.

\textbf{11.} ثلاث أوراق (آس وورقتان أخريان)، وصديق
يسترق النظر ويتخلص دائمًا من ورقة غير الآس من بين الورقتين اللتين لم
تخترهما؛ العب $300$ جولة مبدّلًا دائمًا، وسجّل الفوز.
\omterm{def:g10:stats:series}{والتواتر} المنتظر $\frac23 \approx 0.667$، مع تذبذب
$\pm\frac{1}{\sqrt{300}} \approx 0.058$: أي توقّع ما بين نحو
$61\,\%$ و $72\,\%$ من الفوز بالتبديل --- وهو بعيد عن أي
``نصف بالنصف''.

\textbf{12.} التشكيلات المتساوية \omterm{def:g10:proba:distribution}{الاحتمال} (الأكبر ثم الأصغر):
ولد-ولد، وولد-بنت، وبنت-ولد، وبنت-بنت. و``أحدهما على الأقل ولد'' تبقي ثلاثًا منها:
$P(\text{BB}) = \frac13$.

\textbf{13.} أما ``الأكبر ولد'' فتبقي حالتين:
$P(\text{BB}) = \frac12$. فالمعلومات المختلفة ترشّح
مجموعات سيناريوهات مختلفة --- تمامًا كما فعلت صيغتا مونتي: فما
تتعلمه يهم، و\emph{كيف} تعلمته يهم بالقدر
نفسه.

\textbf{14.} كلها متماثلة: صورة-صورة-صورة، وكتابة-كتابة-كتابة: $\frac28 = \frac14$.
وخطأ دالمبير: أن ``حالتيه'' ليستا متساويتي \omterm{def:g10:proba:distribution}{الاحتمال}
--- إذ إن ``ليست كلها متماثلة'' تحزم ستًّا من الإمكانيات الثماني.
و \Cref{prop:g10:proba:uniform} لا تعدّ إلا على إمكانيات متساوية
\omterm{def:g10:proba:distribution}{الاحتمال}.

\textbf{15.} المكاسب المتوسطة لكل لعبة: $6 \times \frac16 + 3
\times \frac16 = 1.50$ يورو، مقابل رهان $2$ يورو: أي
خسارة متوسطة قدرها $0.50$ لكل لعبة، ونحو $300$ يورو على \omterm{def:g10:stats:quartiles}{مدى} $600$
لعبة. والقسط العادل يساوي المدفوع \omterm{def:g10:stats:mean}{المتوسط}:
$1\,\% \times 10\,000 = 100$ يورو؛ أما الأقساط الحقيقية فتضيف التكاليف
وهامش أمان --- إذ يجب أن يصمد المؤمِّن في السنوات السيئة،
لا في المتوسطة فحسب.

\textbf{16.} \omterm{def:g10:numbers:interval}{الفترة}: $0.5 \pm \frac{1}{\sqrt{100}} =
\intcc{0.4}{0.6}$. فالنتيجة $55$ صورة ($0.55$): داخلها، ولا شيء لافت.
أما $70$ صورة ($0.70$): فبعيدة خارجها --- فشكّك في القطعة النقدية.

\textbf{17.} $\frac{1}{\sqrt{1000}} \approx 0.032$: أي أن مرشحًا نسبته
$50\,\%$ عادلة يحصل عادةً على
$\intcc{0.468}{0.532}$ --- والقيمة $0.52$ تقع مرتاحة داخلها.
فالاستطلاع \emph{يوحي} بتقدّم لكنه لا يبرهن على شيء: فالهامش
يبتلعه.

\textbf{18.} \omterm{def:g10:proba:sample}{العينة} الأولى: \omterm{def:g10:numbers:interval}{الفترة} $0.04 \pm \frac{1}{20} =
\intcc{0}{0.09}$؛ والقيمة المرصودة $0.07$ داخلها: فلا حكم.
\omterm{def:g10:proba:sample}{والعينة} الثانية: $0.04 \pm \frac{1}{50} = \intcc{0.02}{0.06}$؛ والقيمة
المرصودة $0.07$ تقع خارجها: فالادعاء الآن في شك
جدي. \omterm{def:g10:stats:series}{التواتر} نفسه وأداة أحدّ: فالدقة تنمو مثل
$\sqrt n$.

\textbf{19.} ينكمش التذبذب مثل $\frac{1}{\sqrt n}$:
فتربيع حجم \omterm{def:g10:proba:sample}{العينة} أربع مرات ينصّف الضجيج. أما التحيّز --- أي سؤال
الجمهور الخطأ --- فيصمد أمام أي $n$: فمليونا بطاقة مختارة اختيارًا
سيئًا بقيت خاطئة (مسألة نهاية الأسبوع عن كيف تكذب المعطيات في الكتاب السابق).

\textbf{20.} (1) تحقق من تساوي \omterm{def:g10:proba:distribution}{احتمال} الإمكانيات قبل
العدّ --- وهو ما نسيه دالمبير. (2) الأشجار تضرب على طول
الفروع؛ و``واحدة على الأقل'' تستسلم \omterm{def:g10:proba:operations}{للمتممة}.
(3) المعلومة ترشّح السيناريوهات، و\emph{مصدرها}
يغيّر الترشيح --- مونتي العليم في مواجهة مونتي الساقط، و
الأكبر ولد في مواجهة أحدهما ولد. (4) تقترب التواترات من
الاحتمالات بسرعة $\frac{1}{\sqrt n}$، فتحسم محاكاة
صبورة أي جدال --- بما فيه جدال مع عشرة
آلاف كاتب رسالة غاضب.
\end{solution}
